\subsection{Geometric vs.\ Arithmetic Mean: The Core Difference}

The single difference between Huntington-Hill and Webster is the
rounding threshold used for fractional quotas in $(s, s+1)$:
\begin{align*}
\text{HH threshold:}      &\quad \sqrt{s(s+1)} \\
\text{Webster threshold:} &\quad s + 0.5.
\end{align*}
The arithmetic mean of $s$ and $s+1$ is $(s + s+1)/2 = s + 0.5$.
The geometric mean of $s$ and $s+1$ is $\sqrt{s(s+1)}$.
By the AM-GM inequality, $\text{AM} \geq \text{GM}$, with equality
only when $s = s+1$ (impossible), so:
\[
  \sqrt{s(s+1)} < s + 0.5 \quad \text{for all } s \geq 1.
\]

This means HH's threshold is strictly lower than Webster's.
A state with exact quota $q_i$ satisfying
$\sqrt{s(s+1)} < q_i < s + 0.5$
will receive $s+1$ seats under HH but only $s$ seats under Webster.
The set of such states is the locus of HH-Webster disagreement.

For 2020 Census data:
\begin{itemize}
  \item The natural divisor is $d \approx 761{,}169$.
  \item A state in the HH-Webster disagreement region at seat count $s$
    needs population in the interval
    $(761{,}169 \cdot \sqrt{s(s+1)},\; 761{,}169 \cdot (s+0.5))$.
  \item For $s = 1$: this interval is approximately $(1{,}076{,}607, 1{,}141{,}754)$.
\end{itemize}

Montana's 2020 population (1{,}084{,}225) falls in this interval for $s=1$:
\begin{itemize}
  \item $761{,}169 \times \sqrt{1 \times 2} = 761{,}169 \times 1.414 = 1{,}076{,}413$
  \item $761{,}169 \times 1.5 = 1{,}141{,}754$
  \item Montana population: 1{,}084{,}225 --- within the interval.
\end{itemize}
Montana receives 2 seats under Huntington-Hill and would receive 1 seat
under Webster if the divisor were adjusted to produce 435 total seats
under Webster's rounding rule.
(In the 2020 census, the adjusted Webster divisor still gives Montana 2 seats
because other near-threshold states balance out; the specific state that
would differ requires computing the Webster divisor precisely.)

\subsection{Which Method Is ``Better''?}

The choice between HH and Webster is not a mathematical question with
a unique answer but a normative question about which fairness criterion
to prioritise.
\begin{itemize}
  \item HH's relative deviation criterion (minimising the
    max relative difference in per-capita representation) is appropriate
    if one views representational fairness as inherently relative ---
    the ratio between large and small states' per-capita representation
    matters, not the absolute difference.
  \item Webster's least-squares criterion is appropriate if one views
    fairness as minimising aggregate squared error from proportionality ---
    a standard statistical notion of ``closest to proportional.''
\end{itemize}
The Supreme Court confirmed in \textit{Montana} that Congress has
discretion to choose either criterion \citep{montana1992}.
Most academic analysts view Webster as slightly preferable on grounds
of neutrality \citep{balinski1982fair}, while the 1941 Congress
chose HH on grounds of the relative deviation argument.
